\documentclass[10pt,a4paper]{amsart}
\usepackage[margin=19mm]{geometry}
\usepackage{amsmath,amssymb,mathtools}
\usepackage[scr=rsfs,cal=euler]{mathalfa}
\usepackage{xfrac}
\usepackage[unicode,hidelinks,bookmarks=false]{hyperref}
\newcommand{\ie}{i.e.\ }
\newcommand{\cf}{cf.\ }
\newcommand{\resp}{respectively\ }
\newcommand{\Subsection}{\subsection}
\providecommand{\nobreakdash}{\nobreak\hbox{-}}
\setlength{\emergencystretch}{2em}
\multlinegap0pt
\usepackage{bm}
\usepackage{bbm}
\usepackage{dsfont}
\usepackage{xspace}
\usepackage{cancel}
\usepackage{mathtools}
\usepackage{amsthm}
\makeatletter
\@ifundefined{rmk}{\newtheorem{rmk}{Remark}}{}
\makeatother
\title[Nagell-Lutz Theorem for Imaginary Quadratic Fields and class]{Nagell-Lutz Theorem for Imaginary Quadratic Fields and class groups of quadratic fields}
\author{Leena Mondal, Amrutha Chalil, Kalyan Banerjee}
\date{Source: arXiv:2509.07524v2 (CC BY 4.0)}
\begin{document}
\maketitle

\noindent\emph{Source and licence.} Nagell-Lutz Theorem for Imaginary Quadratic Fields and class groups of quadratic fields. Version arXiv:2509.07524v2 (CC BY 4.0), distributed under the Creative Commons Attribution 4.0 International licence, as recorded in the arXiv licence metadata for this exact version and shown on the abstract page https://arxiv.org/abs/2509.07524v2. Full rights notice: licenses/arxiv-2509.07524-CC-BY-4.0.txt.\par\smallskip

\noindent\textit{Editorial scope.} Counted unit: the Rmk whose source label is \texttt{Elliptic curve1} in the article (source0.tex lines 154--179, source label \texttt{Elliptic curve1}), delivered together with the article's own complete original proof of it (source0.tex lines 422--490, one \texttt{proof} environment of 69 lines). No mathematics is written, repaired or re-derived: statement and proof are the article's own text line for line. The proof is detached from the statement in the source: the article states the result at lines 154--179 and prints its proof at lines 422--490, and the proof environment's own header names the statement, so the association is the article's own. Required context retained uncounted: none. Every definition, display and label of those lines is retained unchanged. Omitted from this package and not redistributed: 1--153, 180--421, 491--1036. Nothing inside a retained excerpt is omitted or altered: every retained body is delivered exactly as the article's lines are written, so no omission receipt of any kind accompanies this package. Citations: the retained excerpts carry no citation command at all, so no bibliography entry of the article is transcribed and no reference is redistributed. Reference adaptation: none was needed and none was made; every reference inside the delivered unit resolves inside it, so no unresolved reference or citation is produced. Printed slips kept literal: no printed word, symbol, hypothesis or number of the counted statement or of its proof was altered. Layout and class: the article's own class is \texttt{amsart} with packages including breqn, xcolor, babel, graphicx, framed, ulem, amsmath, amsthm, amssymb, amsfonts, enumerate, inputenc, geometry, comment; the wrapper is the lane's pinned \texttt{amsart} editorial wrapper at 10pt on A4, which restates the theorem-like environments the retained text needs and adds the article's own macro definitions that the retained text needs (none), with extra wrapper packages (bm, bbm, dsfont, xspace, cancel, mathtools). The delivered numbering is the unit's own. Assets: no retained span contains a figure environment, a graphics-inclusion command, a PDF-page inclusion, a table or a TikZ picture; no image file is copied into this package, the package declares no asset dependency and no image file is redistributed with it. Licence: version arXiv:2509.07524v2 of this article is distributed under the Creative Commons Attribution 4.0 International licence, as recorded in the arXiv licence metadata for this exact version and shown on the abstract page https://arxiv.org/abs/2509.07524v2. A full rights notice is delivered at licenses/arxiv-2509.07524-CC-BY-4.0.txt. Characters: every retained excerpt is pure ASCII, so the delivered engine, XeLaTeX with its default Latin Modern text font, reports no missing character.
 \begin{rmk}
     Let \begin{equation}\label{Elliptic curve1}
         E:y^2=x^3+Ax+B ~~with~~ A,B\in \mathbb Q(\sqrt{D}).
     \end{equation}\\ So,
     \begin{align*}
         A&=\frac{A_1+\sqrt{D}A_2}{A_3+\sqrt{D}A_4},\qquad A_1,A_2,A_3,A_4\in{\mathcal{O}_K}\\
         &=\frac{(A_1A_3+DA_2A_4)+\sqrt{D}(A_2A_3-A_1A_4)}{A_3^2-DA_4^2}\\
         &=\frac{A_1^{'}}{A_2{'}}
     \end{align*}
     where $A_1^{'}=(A_1A_3+DA_2A_4)+\sqrt{D}(A_2A_3-A_1A_4)$ and $A_2^{'}=A_3^2-DA_4^2$ and $A_1^{'}\in {\mathcal{O}}_K, A_{2}^{'}\in \mathbb Z$.
\smallskip
Similarly, $B=\frac{B_{1}^{'}}{B_{2}^{'}}$, $B_{1}^{'}\in{\mathcal{O}_K},B_{2}^{'}\in\mathbb Z$.\\
So \eqref{Elliptic curve1} becomes
\begin{align*}
    y^{2}=x^{3}+\frac{A_{1}^{'}}{A_{2}^{'}}x+\frac{B_{1}^{'}}{B_{2}^{'}}.
\end{align*}\label{ellptic curve2}
Let $D=A_{2}^{'}B_{2}^{'}$, multiplying both sides of \eqref{ellptic curve2} by $D^6$, we will get
\begin{align*}
    (D^3y)^2=(D^2x)^3+D^3A_{1}^{'}B_{2}^{'}(D^2x)+D^5A_{2}^{'}B_{1}^{'}.
\end{align*}\label{ellipticurve3}
Again, let $X=D^2x, Y=D^3y, A^{'}=D^3A_{1}^{'}B_{2}^{'}, B^{'}=D^5A_{2}^{'}B_{1}^{'}$ and then \eqref{ellipticurve3} becomes\\
\begin{equation}
    Y^{2}=X^3+A^{'}X+B^{'}, \qquad A^{'},B^{'}\in{\mathcal{O}_K}.
\end{equation}
Thus, we can assume that our elliptic curve \eqref{Elliptic curve1} has coefficients from $\mathcal{O}_K$.
 \end{rmk}

\begin{proof}
    Since $x,y\in \mathbb Z(\frac{1+\sqrt{D}}{2})$, we can take $x=a+b(\frac{1+\sqrt{D}}{2})$ and $y=c+d(\frac{1+\sqrt{D}}{2})$. That is,\\
    our \textbf{claim} will be $b$ and $d$ are even.\\
    We have from \eqref{Elliptic curve1},\\
    \begin{align*}
        y^2=x^3+Ax+B,~~~A=A_1+A_2\left(\frac{1+\sqrt{D}}{2}\right), B=B_1+B_2\left(\frac{1+\sqrt{D}}{2}\right),
    \end{align*}
    so,
    \begin{align*}
        \left(c+d(\frac{1+\sqrt{D}}{2})\right)^2&=\left(a+b\left(\frac{1+\sqrt{D}}{2}\right)\right)^3+A\left(a+b(\frac{1+\sqrt{D}}{2})\right)+B\\
        &=\left(a+b\left(\frac{1+\sqrt{D}}{2}\right)\right)^3+\left(A_1+A_2(\frac{1+\sqrt{D}}{2})\right)\left(a+b\left(\frac{1+\sqrt{D}}{2}\right)\right)\\&+\left(B_1+B_2\left(\frac{1+\sqrt{D}}{2}\right)\right).
    \end{align*}
    After expanding everything, we will get\\
    \begin{align*}
        c^2+\frac{d^2}{4}+\frac{d^2D}{4}+\frac{d^2\sqrt{D}}{2}+cd+cd\sqrt{D}&=a^3+\frac{a^2b}{2}+\frac{a^2b\sqrt{D}}{2}+\frac{ab^2}{4}+\frac{ab^2\sqrt{D}}{2}+\frac{ab^2D}{4}+\frac{b^3}{8}+\\&\frac{b^3\sqrt{D}}{4}+\frac{b^3D}{8}+\frac{b^3\sqrt{D}}{8}+\frac{b^3D}{4}+\frac{b^3D^{3/2}}{8}+A_1a+\frac{A_1b}{2}+\\&\frac{A_1b\sqrt{D}}{2}+\frac{A_2a}{2}+\frac{A_2a\sqrt{D}}{2}+\frac{A_2b}{4}+\frac{A_2bD}{4}+\frac{A_2bD^4}{2}.
    \end{align*}
    By comparing the real part, we will get
    \begin{align*}
        c^2+\frac{d^2}{4}+\frac{d^2D}{4}+cd&=a^3+\frac{a^2b}{2}+\frac{ab^2}{4}+\frac{ab^2D}{4}+\frac{b^3}{8}+\frac{b^3D}{4}+A_1a+\frac{A
        _1b}{2}+\frac{A_2a}{2}\\&+\frac{A_2b}{4}+\frac{A_2bD}{4}+\frac{A_2bD}{2},
    \end{align*}
    which will give,\\
    \begin{align*}
        2d^2+2d^2D\equiv 2ab^2+2ab^2D+b^3+b^3D+2b^3D+2A_2b+2A_2bD~(mod~4),
    \end{align*}
    $\implies$
    \begin{align*}
         -4d^2\equiv -4ab^2-2b^3-4A_2b~(mod~4),
    \end{align*}
    $\implies$
    \begin{align*}
        -2b^3\equiv0~(mod~4)
    \end{align*}
   $\implies b^3\equiv0~(mod~4)\implies b=0,2~(mod~4)$.\\ 
   Then $b=4k$ or $b=4k+2$.\\
    And by comparing the imaginary part, we get
    \begin{align*}
        \frac{d^2}{2}+cd=\frac{a^2b}{2}+\frac{ab^2}{2}+\frac{b^3}{8}+\frac{b^3D}{8}+\frac{A_1b}{2}+\frac{A_2a}{2},
    \end{align*}
    which gives\\
    \begin{align*}
        4d^2+8cd=4a^2b+4ab^2+3b^3+b^3D+4A_1b+4A_2a.
    \end{align*}
    For $b=4k$, we get
    \begin{align*}
        d^2+2cd=4a^2k^2+4^2k^2a+3\times4^2k^3+4^2k^3D+4A_1k+A_2a
    \end{align*}
    gives,
    \begin{equation}\label{A2}
        d^2\equiv A_2a~(mod~2).
    \end{equation}
    Since the discriminant of \eqref{Elliptic curve1} is non-zero, we have
    \begin{align*}
        4\left(A_1+A_2\left(\frac{1+\sqrt{D}}{2}\right)\right)^3+27\left(B_1+B_2\left(\frac{1+\sqrt{D}}{2}\right)\right)^2\neq 0
    \end{align*}
    \begin{align*}
       0&\neq 4A_1^3+\frac{A_1^2A_2}{2}+\frac{A_1^2A_2\sqrt{D}}{2}+\frac{A_1A_2^2}{2}+\frac{A_1A_2^2\sqrt{D}}{2}+\frac{A_2^3+DA_2^3+2A_2^3\sqrt{D}}{4}\\&+A_1^2A_2+A_1^2A_2\sqrt{D}+\frac{A_1A_2^2}{2}+\frac{DA_1A_2^2}{2}+A_1A_2^2\sqrt{D}+27B_1^2+\frac{27B_2^2}{4}+\frac{27DB_2^2}{4}\\&+\frac{27B_2^2\sqrt{D}}{2}B_1B_2+B_1B_2\sqrt{D},
    \end{align*}
    which gives,
    \begin{align*}
        2A_1^2A_2+2A_1A_2^2+A_3^2+2A_1A_2^2-6A_1A_2^2+3B_2^2-9DB_2^2\equiv 0~(mod~4),
    \end{align*}
    so,
    \begin{align*}
        2A_1^2A_2-4A_1A_2^2+A_2^3+2A_1A_2^2-6B_2^2\equiv0~(mod~2).
    \end{align*}
    $\implies A_2^3\equiv 0~(mod~2)\implies A_2\equiv 0~(mod~2)\implies A_2$ is even. \\
    Therefore, from \eqref{A2} we can conclude that $d$ is even. One can proceed the same way for $b=4k+2$. Hence our claim and we have the following theorem.
\end{proof}

\end{document}
